% -*- coding: utf-8 -*-
% !TEX program = xelatex
\documentclass[plain,noanswer,medmath]{../../template/jnuexam}
\usepackage{../../template/kysx}

\begin{document}


\examtitle{2001}{2}


\loadproblems{1}{2001P1}%载入试卷一


\makepart{填空题}{本题共5小题，每小题3分，满分15分}


\begin{problem}
$\lim_{x \to 1} \frac{\sqrt{3 - x} - \sqrt{1 + x}}{x^2 + x - 2} =$ \fillin{}．
\end{problem}


\begin{problem}
设函数$y = f(x)$由方程$\e^{2x + y} - \cos(xy) = \e - 1$所确定，则曲线$y = f(x)$在点$(0,1)$处的法线方程为 \fillin{}．
\end{problem}


\begin{problem}
$\int_{- \frac{\pi}{2}}^{\frac{\pi}{2}} \left(x^3 + \sin^2 x \right) \cos^2 x\dx =$ \fillin{}．
\end{problem}


\begin{problem}
过点$\left(\frac{1}{2},0 \right)$且满足关系式$y'\arcsin x + \frac{y}{\sqrt{1 - x^2}} = 1$的曲线方程为 \fillin{}．
\end{problem}


\begin{problem}
设方程$\begin{pmatrix}
  a & 1 & 1 \\
  1 & a & 1 \\
  1 & 1 & a
\end{pmatrix}\begin{pmatrix}
  x_1 \\
  x_2 \\
  x_3
\end{pmatrix} = \begin{pmatrix}
  1 \\
  1 \\
  -2
\end{pmatrix}$有无穷多个解，则$a =$ \fillin{}．
\end{problem}


\makepart{选择题}{本题共5小题，每小题3分，共15分}


\begin{problem}
设$f(x) = \begin{cases}
  1, & |x| \le 1, \\
  0, & |x| > 1,
\end{cases}$则$f\left\{f\left[f(x) \right] \right\}$等于\pickout{}
\begin{abcd}
\item $0$．
\item $1$．
\item $\begin{cases}
  1, & |x| \le 1, \\
  0, & |x| > 1.
\end{cases}$
\item $f(x) = \begin{cases}
  0, & |x| \le 1, \\
  1, & |x| > 1.
\end{cases}$
\end{abcd}
\end{problem}


\begin{problem}
设当$x \to 0$时，$(1 - \cos x)\ln(1 + x^2)$是比$x\sin x^n$高阶的无穷小，
$x\sin x^n$是比$\left(\e^{x^2} - 1 \right)$高阶的无穷小，则正整数$n$等于\pickout{}
\begin{abcd}
\item $1$．
\item $2$．
\item $3$．
\item $4$．
\end{abcd}
\end{problem}


\begin{problem}
曲线$y = (x - 1)^2(x - 3)^2$的拐点个数为\pickout{}
\begin{abcd}
\item $0$．
\item $1$．
\item $2$．
\item $3$．
\end{abcd}
\end{problem}


\begin{problem}
已知函数$f(x)$在区间$(1 - \delta ,1 + \delta)$内具有二阶导数，$f'(x)$严格单调减少，且$f(1) = f'(1) = 1$，则\pickout{}
\begin{abcd}
\item 在$(1 - \delta ,1)$和$(1,1 + \delta)$内均有$f(x) < x$．
\item 在$(1 - \delta ,1)$和$(1,1 + \delta)$内均有$f(x) > x$．
\item 在$(1 - \delta ,1)$内$f(x) < x$，在$(1,1 + \delta)$内$f(x) > x$．
\item 在$(1 - \delta ,1)$内$f(x) > x$，在$(1,1 + \delta)$内$f(x) < x$．
\end{abcd}
\end{problem}


\useproblem{1}{2}{1}%【同试卷一第二(1)题】


\makepart{}{本题满分6分}

\begin{problem}
求$\int \frac{\dx}{(2x^2 + 1)\sqrt{x^2 + 1}}$．
\end{problem}


\makepart{}{本题满分7分}

\begin{problem}
求极限$\lim_{t \to x} \left(\frac{\sin t}{\sin x} \right)^{\frac{x}{\sin t - \sin x}}$，
记此极限为$f(x)$，求函数$f(x)$的间断点并指出其类型．
\end{problem}


\makepart{}{本题满分7分}

\begin{problem}
设$\rho = \rho(x)$是抛物线$y = \sqrt{x}$上任一点$M(x,y)$（$x \ge 1$）处的曲率半径，
$s = s(x)$是该抛物线上介于点$A(1,1)$与$M$之间的弧长，
计算$3\rho \frac{\d^2\rho}{\ds^2} - \left(\frac{\d\rho}{\ds} \right)^2$的值%
（在直角坐标系下曲率公式为$K = \frac{\left| y'' \right|}{(1 + y'^2)^{3/2}}$）．
\end{problem}


\makepart{}{本题满分7分}

\begin{problem}
设函数$f(x)$在$[0, +\infty)$上可导，$f(0) = 0$，且其反函数为$g(x)$．若
$$\int_0^{f(x)} g(t)\dt = x^2\e^x,$$
求$f(x)$．
\end{problem}


\makepart{}{本题满分7分}

\begin{problem}
设函数$f(x),g(x)$满足$f'(x) = g(x),g'(x) = 2\e^x - f(x)$，且$f(0) = 0,g(0) = 2$，求
$$\int_0^{\pi}\left[\frac{g(x)}{1 + x} - \frac{f(x)}{(1 + x)^2} \right] \dx.$$
\end{problem}


\makepart{}{本题满分9分}

\begin{problem}
设$L$是一条平面曲线，其上任意一点$P(x,y)$（$x > 0$）到坐标原点的距离恒等于该点处的切线在$y$轴上的的截距，
且$L$经过点$\left(\frac{1}{2},0 \right)$．
\begin{enumerate}
\item 试求曲线$L$的方程．
\item 求$L$位于第一象限部分的一条切线，使该切线与$L$以及两坐标轴所围图形面积最小．
\end{enumerate}
\end{problem}


\makepart{}{本题满分7分}

\begin{problem}
一个半球体状的雪堆，其体积融化的速率与半球面面积$S$成正比，比例常数$K > 0$．
假设在融化过程中雪堆始终保持半球体状，已知半径为$r_0$的雪堆在开始融化的$3$小时内，
融化了其体积的$\frac{7}{8}$，问雪堆全部融化需要多少小时？
\end{problem}


\makepart{}{本题满分8分}

\begin{problem}
设$f(x)$在区间$[- a,a]$（$a > 0$）上具有二阶连续导数，$f(0) = 0$．
\begin{enumerate}
\item 写出$f(x)$的带拉格朗日余项的一阶麦克劳林公式；
\item 证明在$[- a,a]$上至少存在一点$\eta$，使$a^3 f''(\eta) = 3\int_{- a}^a f(x)\dx$．
\end{enumerate}
\end{problem}


\makepart{}{本题满分6分}

\begin{problem}
已知矩阵$A = \begin{pmatrix}
  1 & 0 & 0 \\
  1 & 1 & 0 \\
  1 & 1 & 1
\end{pmatrix}$，$B = \begin{pmatrix}
  0 & 1 & 1 \\
  1 & 0 & 1 \\
  1 & 1 & 0
\end{pmatrix}$．且矩阵$X$满足
$$AXA + BXB = AXB + BXA + E,$$
其中$E$是$3$阶单位矩阵，求$X$．
\end{problem}


\makepart{}{本题满分6分}%【类似试卷一第九题】

\begin{problem}
设$\alpha_1,\alpha_2,\alpha_3,\alpha_4$为线性方程组$AX = 0$的一个基础解系，
$$\beta_1 = \alpha_1 + t\alpha_2,\beta_2 = \alpha_2 + t\alpha_3,\beta_3 = \alpha_3 + t\alpha_4,\beta_4 = \alpha_4 + t\alpha_1.$$
试问实数$t$满足什么关系时，$\beta_1,\beta_2,\beta_3,\beta_4$也为$AX = 0$的一个基础解系．
\end{problem}


\end{document}
